We prove the Gaussian certificate, the outer optimization, and the normalized primitive transfer in turn.

\emph{Step 1: unbounded Gaussian certification.}
Fix a terminal covariance cell and write \(d=|A|\).  Compactness and uniform positive definiteness give a computable bound \(\bar\sigma\) on all coordinate standard deviations.  The rule selecting the largest Gaussian coordinate has, for a true best arm \(b\) and gaps \(d_i=h_b-h_i\), risk at most
\[
\sum_{i\ne b}d_i\Phi\!\left(-\frac{d_i}{\sqrt{\sigma_i^2+\sigma_b^2}}\right)
\leq\frac{(d-1)\bar\sigma}{\sqrt\pi}=:H_0. \tag{51}
\]
The last inequality follows from the Mills calculation in \(\texttt{lem:hierarchical-gaussian-screening}\) and \(\sqrt{\sigma_i^2+\sigma_b^2}\leq\sqrt2\bar\sigma\).  Hence \(G_A\leq H_0\) throughout the cell.

Choose rational \(0<\lambda<1\) so that
\[
\{(1-\lambda)^{-1/2}-1\}H_0\leq\varepsilon/4, \tag{52}
\]
then rational \(T>0\) so that, for \(a_\lambda=\sqrt{2/\lambda}\bar\sigma\),
\[
\frac{2d^2a_\lambda}{\sqrt{2\pi}}
\exp\!\left\{-\frac{T^2}{2a_\lambda^2}\right\}\leq\varepsilon/4. \tag{53}
\]
The searches terminate because the left side of (52) tends to zero as \(\lambda\downarrow0\), and that of (53) tends to zero as \(T\to\infty\).  Directed rational interval arithmetic certifies both comparisons.  For every nonempty \(S\subseteq A\), choose rational
\[
M_S\geq2\sqrt{1-\lambda}\,T\sqrt{|S|-1}. \tag{54}
\]
For a singleton use its exact zero value.  For every larger \(S\), invoke \(\texttt{thm:gaussian-finite-certificate}\) with gap \(\eta\) satisfying \((1-\lambda)^{-1/2}\eta\leq\varepsilon/4\), and write \(L_S\leq G_{S,M_S}\leq U_S\), \(U_S-L_S\leq\eta\).  Define
\[
L^{\infty}_{\varepsilon,A}=\max_{\varnothing\ne S\subseteq A}L_S,
\quad
U^{\infty}_{\varepsilon,A}=(1-\lambda)^{-1/2}
\max_{\varnothing\ne S\subseteq A}U_S
+\frac{2d^2a_\lambda}{\sqrt{2\pi}}e^{-T^2/(2a_\lambda^2)}. \tag{55}
\]
The hierarchical screening lemma gives \(L^{\infty}_{\varepsilon,A}\leq G_A\) and the certified selector upper bound.  If \(S_\star\) maximizes \(U_S\), then
\[
\max_SU_S=U_{S_\star}\leq L_{S_\star}+\eta\leq\max_SL_S+\eta. \tag{56}
\]
Since \(\max_SL_S\leq G_A\leq H_0\), (52)--(56) imply
\[
U^{\infty}_{\varepsilon,A}-L^{\infty}_{\varepsilon,A}
\leq\{(1-\lambda)^{-1/2}-1\}H_0
+(1-\lambda)^{-1/2}\eta
+\frac{2d^2a_\lambda}{\sqrt{2\pi}}e^{-T^2/(2a_\lambda^2)}
<\varepsilon. \tag{57}
\]
There are only \(2^d-1\) retained subsets.  The receipt records \(\lambda,T,M_S\), every compact checksum, the split, tail enclosure, and maxima in (55).  No radial bound on \(h\) was used.

\emph{Step 2: outer-simplex certification.}
For Bernoulli labels, positivity of the binomial probabilities gives
\[
(Bp)_k\geq\underline q_k:=\min_{\ell\in\mathcal K}B_{k\ell}>0
\quad(p\in\Delta_K,\ k\in A). \tag{58}
\]
Partition the explicit constraints \(p_\ell\geq0\), \(\sum_\ell p_\ell=1\) into rational simplex cells.  Linear interval evaluation supplies \(q^-\leq Bp\leq q^+\).  By \(\texttt{lem:gaussian-information-order}\), evaluation at \(q^+\) is a lower bound for the cell, while evaluation at a feasible rational \(p\) is an upper incumbent.  If the relative coordinate width is at most \(\xi<1\), the same lemma bounds unresolved multiplicative variation by
\[
(1-\xi)^{-1/2}-(1+\xi)^{-1/2}. \tag{59}
\]
The lower bound (58) makes (59) tend uniformly to zero under refinement.  Discard every branch whose lower value exceeds the incumbent and refine until the incumbent-minus-best-lower gap is at most \(\varepsilon\).  Compactness ensures termination; every minimizer remains in the surviving rational cells.

For exact counts, optimize over
\[
a_i\geq0\quad(i\in A),\qquad\sum_{i\in A}a_i=1. \tag{60}
\]
At \(a_i=0\), \(\texttt{lem:gaussian-information-order}\) gives \(+\infty\).  In the positive interior, the subface bound and \(\texttt{thm:two-policy-minimax}\) give, for distinct \(i,j\),
\[
G_A(v_A,a)\geq G_{\{i,j\}}(v_{\{i,j\}},a_{\{i,j\}})
=\kappa\sqrt{\frac{v_i}{a_i}+\frac{v_j}{a_j}}
\geq\kappa\sqrt{\frac{v_i}{a_i}}. \tag{61}
\]
Let \(\underline v>0\) be the cellwise lower bound and let \(\overline U\) be a certified finite upper bound at the uniform allocation.  Choose rational
\[
0<\eta_0<\min\left\{\frac1{2d},
\frac{\kappa^2\underline v}{4(\overline U+1)^2}\right\}. \tag{62}
\]
If any \(a_i<\eta_0\), (61) exceeds \(2(\overline U+1)>\overline U\), so no minimizer lies there.  The optimizer is inside the compact polytope (60) with \(a_i\geq\eta_0\), where the preceding positive-cell branch-and-bound terminates.  This proves both outer computations without evaluating a positive-rate formula at zero.

\emph{Step 3: normalized primitive transfer.}
Use prespecified disjoint cluster indices for a covariance pilot, screen, and decision.  Let
\[
 m_C+N_C^{\mathrm s}+N_C^{\mathrm d}=C,\qquad
m_C\to\infty,\quad m_C/C\to0,\qquad
N_C^{\mathrm s}/C\to\lambda,\qquad
N_C^{\mathrm d}/C\to1-\lambda. \tag{63}
\]
Under complete randomization, make the split separately inside every active stratum with these limiting fractions.  Schedule independence and the respective randomization assumptions make the three outcome arrays independent conditional on their prespecified assignments and counts.

Write \(\widetilde U_C^{\mathrm s}\) and \(\widetilde U_C^{\mathrm d}\) for the cell-calibrated estimators formed only from the screening and decision pieces.  Under Bernoulli labels, every relevant cell has probability
\[
\rho_z=\frac{q_k}{M_k}
\geq\frac{\min_\ell B_{k\ell}}{M_k}>0,
\qquad z\in\mathcal S_k. \tag{64}
\]
Under exact counts on the certified interior \(a_k\geq\eta_0\), its asymptotic frequency is at least \(\eta_0/M_k\).  Since \(n,K\) are fixed, the reciprocal cell probabilities are uniformly bounded.

Expanding each cell mean around its deterministic cell frequency gives, for either piece \(b\in\{\mathrm s,\mathrm d\}\),
\[
\widetilde U_{C,k}^{b}-U_k
=\frac1{N_C^b}\sum_{c\in b}
\frac{\mathbf1\{Z_c\in\mathcal S_k\}}{q_k}
\{W_c(Z_c)-\mu_{Z_c}\}+o_P((N_C^b)^{-1/2}) \tag{65}
\]
under Bernoulli labels, with \(q_k\) replaced by the active exact-count share under complete randomization.  The finitely many ratio remainders are uniform because of (64) and bounded outcomes.  Bounded characteristic-function expansions, exactly as in the two compact oracle transfers, now give jointly
\[
\sqrt C(\widetilde U_C^{\mathrm s}-U)
\Rightarrow N(0,\Sigma/\lambda),\qquad
\sqrt C(\widetilde U_C^{\mathrm d}-U)
\Rightarrow N(0,\Sigma/(1-\lambda)), \tag{66}
\]
and the limits are independent.  Thus the primitive screen uses differences of \(\sqrt C\widetilde U_C^{\mathrm s}\), while the compact decision rule receives
\[
\sqrt{1-\lambda}\,\sqrt C\,\widetilde U_C^{\mathrm d}. \tag{67}
\]
Translation invariance removes its unknown common location.  On a local welfare displacement \(h\), (67) has contrast mean \(\sqrt{1-\lambda}\,h\) and covariance \(\Sigma\), exactly matching the Gaussian construction and its regret rescaling in the screening lemma.  Equations (66)--(67) are the required normalization; using the unscaled decision statistic would be wrong.

The pilot empirical second moments are uniformly consistent by boundedness and (64), so outward covariance snapping selects the certified cell or a recorded neighbor with probability tending to one.  A binomial lower-tail bound for cell counts and Hoeffding's inequality conditional on assignments give constants \(c_1,c_2>0\), uniform on a positive design cell, such that
\[
\Pr\{\sqrt C\,|\widetilde U_{C,k}^{\mathrm s}-U_k|>x\}
\leq c_1e^{-c_2x^2}+c_1e^{-c_2C},
\qquad0\leq x\leq\sqrt C. \tag{68}
\]
Repeating the pairwise deletion and bad-range decomposition of \(\texttt{lem:hierarchical-gaussian-screening}\) with (68) gives a deterministic \(b(T)\downarrow0\) such that the expected scaled regret from deleting a best arm or retaining a set of scaled range above \(2T\) is at most \(b(T)+o(1)\).  In particular, a coordinate whose scaled welfare gap diverges contributes \(o(1)\).

To prove uniformity, suppose a violating triangular sequence existed.  The schedule space \([0,1]^{n\times\mathcal Z_n}\) is compact metric, so \(\texttt{lem:weak-compactness-compact-laws}\) gives a subsequence on which the schedule laws converge weakly.  Coordinate welfare and its square are bounded continuous functions, so all target cell means and variances converge; compactness of the certified covariance cell supplies the required positive-definite covariance limit.  Pass again so every nonnegative scaled gap converges in \([0,\infty]\).  Infinite-gap coordinates have \(o(1)\) regret by (68).  On the remaining finite-gap coordinates, (65)--(67) give joint pilot-screen-decision Gaussian convergence.  Gaussian contrasts and the finite represented-rule boundaries have probability zero, so the limiting risk is that of the corresponding Gaussian screened selector and is bounded by (55), contradicting the violation.  Therefore the primitive limit superior is bounded uniformly by \(U^{\infty}_{\varepsilon,A}\).

This argument ranges over retained subsets \(S\subseteq A\) only.  With \(A=\mathcal K\), (64) gives positive information for every Bernoulli candidate and hence global adaptation.  Under exact counts, global adaptation additionally requires all candidate shares to remain in a compact interior set.  At a vanishing action-relevant share, only the extended value \(+\infty\) from (61) is asserted.

Finally, for stage \(j\), take tolerance \(1/j\), choose \(\lambda_j,T_j,\eta_j\) by (52)--(54), and then choose \(C_j\) so large that the uniform pilot, split-CLT, snapping, and screening errors are at most \(1/j\), with \(m_C\geq j\), \(m_C/C\leq1/j\), \(N_C^{\mathrm s}\geq j\), and \(N_C^{\mathrm d}\geq j\).  Use that stage for \(C_j\leq C<C_{j+1}\).  Then the pilot fraction vanishes, every stagewise screen and decision sample diverges, and (57) plus the primitive transfer makes the certified risk converge to \(G_A\).  Combining this diagonal schedule with Step 2 proves the asserted certified design and implementation.